% Table 4.1 — External validation against the benchmark's originating work (AEGIS/VLSA).
% Requires: \usepackage{booktabs}  \usepackage{threeparttable}
%
% NUMBERS AND THEIR PROVENANCE:
%   This work      — figures/results_shielded_distillation.md, eval 2026-08-07, n = 120 held-out.
%   AEGIS          — their Table I, averaged over the three suites used here (spatial, goal,
%                    object). Their own averages additionally include a Long suite, so these are
%                    recomputed three-suite figures, NOT their printed averages. Verified against
%                    the PDF 2026-08-18; see papers/FINDINGS.md.
%   Do NOT cite 87.5% for their shielded CAR. That figure came from a docstring in their released
%   code, not the paper, and is wrong. The paper's three-suite figure is 71.9%.

\begin{table}[htbp]
\centering
\begin{threeparttable}
\caption[External validation against the published baseline]{%
  External validation. The unshielded baseline measured here reproduces the published
  $\pi_{0.5}$ figures to within a fraction of a point on both metrics, indicating that the
  evaluation pipeline recovers the published result rather than an artefact of this
  implementation. The shielded rows are \emph{not} like-for-like and are discussed below.}
\label{tab:external_validation}
\begin{tabular}{llcc}
\toprule
\textbf{Condition} & \textbf{Source} & \textbf{CAR} & \textbf{TSR} \\
\midrule
\multirow{2}{*}{Base $\pi_{0.5}$, no shield}
  & Published\tnote{a}      & 17.3\%          & 58.9\%          \\
  & This work\tnote{b}      & \textbf{17.5\%} & \textbf{58.3\%} \\
\addlinespace[2pt]
\multicolumn{2}{l}{\emph{difference}} & $0.2$ pp & $0.6$ pp \\
\midrule
\multirow{2}{*}{With runtime CBF shield}
  & Published\tnote{a}      & 71.9\%          & 74.6\%          \\
  & This work\tnote{b}      & 86.7\%          & 71.7\%          \\
\addlinespace[2pt]
\multicolumn{2}{l}{\emph{difference}\tnote{c}} & $+14.8$ pp & $-2.9$ pp \\
\bottomrule
\end{tabular}
\begin{tablenotes}[flushleft]\footnotesize
\item[a] Averaged over the three suites used in this work (spatial, goal, object) from the
  originating work's reported per-suite results in the full action space. Their printed averages
  additionally include a Long suite and are therefore not directly quoted here.
\item[b] $n = 120$ held-out rollouts, pooled over 24 scenes. CAR is the complement of the collision
  rate defined in Section~\ref{sec:metrics}; note that the originating work defines collision
  avoidance as the proportion of strictly collision-free episodes without stating a numeric
  tolerance, so the 1\,mm threshold used here is this work's operationalisation of that definition.
\item[c] The shielded rows are not a like-for-like comparison and the gap is expected. Two
  differences in the constraint account for it, both deliberate: the published formulation
  constrains only the end-effector, by its authors' own statement, whereas the barrier used here
  additionally constrains links three to seven and the hand (Section~\ref{sec:barrier}); and this
  barrier reads ground-truth geometry from the simulator, whereas theirs is grounded through
  vision-language perception and fused depth, to which they attribute their own residual
  collisions. Both differences favour the shield used here, so these figures are an upper bound
  for this class of filter rather than an improvement over that method.
\end{tablenotes}
\end{threeparttable}
\end{table}
